% sections/frontmatter/abstract.tex
\begin{abstract}
Il sistema di identificazione automatica (AIS) produce ogni giorno grandi quantità di dati sulla
posizione e sul moto delle navi. Questi dati permettono di individuare in modo automatico i
comportamenti anomali nel traffico marittimo. La tesi re-implementa in Python il framework TREAD
di Pallotta et al.\ e lo applica ai messaggi AIS pubblicati dal NOAA per l'arcipelago delle Hawaii
nel primo trimestre del 2019. L'obiettivo è verificare se un metodo sviluppato agli inizi degli
anni 2010 sia ancora applicabile, con gli strumenti attuali, a un dataset reale di milioni di
record.

Il lavoro si sviluppa in tre fasi. Vengono prima puliti i dati e classificati gli stati delle navi, poi
vengono estratti senza supervisione i waypoint e le rotte che li collegano per poi infine assegnare 
a ogni trip (viaggio) di una rotta un punteggio di normalità costruito secondo i criteri stabiliti dall'algoritmo. 
Per gestire circa 5,4 milioni di record senza hardware specializzato l'elaborazione
avviene in modalità batch e i punti stazionari vengono aggregati prima del clustering.

Il metodo estrae waypoint, tra cui i principali porti dell'arcipelago, e rotte coerenti con il traffico reale. 
È stata poi scelta una rotta, in questo caso che collega il porto di Honolulu al porto di Kalaeloa, 
dalla quale sono stati valutati 184 trip svoltisi nel periodo di analisi. I risultati dell'analisi
sono stati confrontati con un'etichettatura manuale e da questo confronto sono state calcolate alcune metriche.
Per questa rotta specifica, con la soglia scelta, non si verifica alcun falso positivo, ma viene rilevato solo un quinto circa delle anomalie.
Le prestazioni migliorano se si considerano solo le anomalie di tipo route deviation, ossia la tipologia
di anomalie per cui questo metodo è stato pensato. In questo caso ne viene rilevata circa metà.
Il limite principale riguarda le zone portuali, dove l'addensamento dei punti agisce sul punteggio di normalità.

La tesi propone infine alcuni sviluppi, tra cui modelli separati per le zone portuali e una versione del metodo non più a batch.
\end{abstract}
